%!TEX root = ../mainufficiale.tex
\chapter{Conclusioni}
\label{cap:conclusioni}

Il Capitolo~\ref{cap:introduzione} ha posto il controllo semaforico su rete
come un problema di coordinamento distribuito, in cui la fase scelta a un
incrocio incide sul traffico che raggiunge gli incroci vicini, e ha
chiesto a un controllore di perseguire insieme due proprietà distinte,
l'efficienza nel far scorrere il traffico e l'equità tra le direzioni che
attendono di essere servite. Questo lavoro affronta il problema
addestrando un agente di reinforcement learning per intersezione,
coordinato con i vicini tramite una rete neurale a grafo, a cui è
aggiunto un termine di ricompensa dedicato all'attesa massima osservata
su una corsia. Il Capitolo~\ref{cap:confronto} ha riportato i risultati
di tre studi sperimentali, condotti sulle stesse dieci configurazioni di
traffico, tre di training e sette mai incontrate in addestramento.
Questo capitolo li riprende alla luce delle tre domande di ricerca poste
nell'introduzione, indica le direzioni lasciate aperte e discute infine i
limiti dello studio.

Un controllore appreso e coordinato tramite reti neurali a grafo riduce
l'attesa massima subita dai veicoli rispetto a Max-Pressure, a un costo
contenuto in termini di efficienza media. Il modello di riferimento di
questo lavoro riduce l'attesa massima del 59\% in training e del 52\% in
test rispetto a Max-Pressure, a fronte di un tempo di percorrenza medio
superiore rispettivamente del 6.7\% e dell'1.1\%, un margine che si
annulla del tutto sulle due griglie $6\times6$, dove il modello registra
anzi un tempo di percorrenza inferiore a quello di Max-Pressure. Lo
stesso modello completa inoltre una quota di veicoli più alta di
Max-Pressure su tutte le sette configurazioni di test. Il confronto di
base rivela già un aspetto meno atteso di questo compromesso:
Max-Pressure non riesce a battere nemmeno Fixed-Time, il piano a
rotazione fissa più semplice tra i controllori confrontati, né sul tempo
di percorrenza massimo né sull'attesa massima, in nessuno dei due regimi:
la sua scelta miope della fase a pressione più alta, senza considerare da
quanto tempo attende la direzione esclusa (Capitolo~\ref{cap:introduzione}),
non pone alcun limite strutturale all'attesa, a differenza della
rotazione cieca ma periodica di Fixed-Time. Solo il modello di questo
lavoro batte entrambe le baseline su ogni metrica, in entrambi i regimi.

Variando il peso assegnato alla componente di equità nella ricompensa,
questo lavoro quantifica direttamente il compromesso appena descritto. Al
diminuire di questo peso, il tempo di percorrenza massimo peggiora in
modo strettamente monotono e l'attesa massima in modo quasi monotono; il
tempo di percorrenza medio non segue invece un andamento altrettanto
prevedibile, tanto che azzerare del tutto la penalità sull'attesa non
produce il tempo di percorrenza più basso tra le varianti confrontate, né
in training né in test. Una spiegazione plausibile è che il segnale di
attesa per corsia resti comunque presente nello stato osservato dal
modello anche quando la ricompensa smette di penalizzarlo, un
disallineamento tra ciò che il modello osserva e ciò per cui viene
premiato che può disturbare l'addestramento. Azzerare del tutto questo
peso non riporta comunque il controllore al comportamento di
Max-Pressure: pur restando la configurazione peggiore tra quelle provate
su questo asse, resta nettamente sotto sia Max-Pressure sia Fixed-Time
sulle metriche di caso peggiore, in entrambi i regimi.

Tale vantaggio si mantiene, con qualche erosione attesa, su topologie di
rete, spaziature stradali e profili di traffico mai incontrati in
addestramento. Il margine di riduzione dell'attesa massima rispetto a
Max-Pressure resta ampio ma si restringe passando dal training al test,
dal 59\% al 52\%; una lettura cauta è che parte di questo vantaggio si
eroda su condizioni mai incontrate in addestramento, senza che questo
permetta di per sé conclusioni più forti sulle capacità di
generalizzazione del modello. Il vantaggio sul completamento del
percorso, al contrario, si mantiene su ogni singola configurazione di
test. Tra le modifiche isolate nello studio di ablazione, solo la
rimozione del meccanismo che permette alla rete ricorrente di propagare
informazione lungo la sequenza temporale mostra un divario che cresce con
la scala della griglia, da circa il 10\% sulle configurazioni $4\times4$
a circa il 12\% sulla $6\times6$, mentre le altre due modifiche che
toccano un singolo dettaglio dell'algoritmo di apprendimento restano
vicine al riferimento senza un vantaggio sistematico su di esso. La sola
modifica che disattiva insieme il segnale di attesa nello stato, la
maschera anti-starvation e il vincolo sul campo visivo degrada invece in
modo pressoché uniforme su tutte le configurazioni fuori distribuzione,
circa il 10\% su quelle a spaziatura diversa e il 12\% sulle altre, senza
che la spaziatura risulti un fattore distintivo.

Tra le scelte progettuali studiate in questo lavoro, quelle relative allo
stato e alla ricompensa risultano più determinanti di quelle relative ai
soli dettagli dell'algoritmo di apprendimento. Le tre varianti che
disattivano singolarmente un meccanismo dell'algoritmo di apprendimento,
la propagazione del gradiente lungo la sequenza temporale, la doppia
stima del valore atteso nell'aggiornamento, o alcuni accorgimenti della
memoria delle esperienze passate, restano vicine al riferimento su quasi
ogni configurazione. Le due varianti che rimuovono invece congiuntamente
il segnale di attesa nello stato e la maschera anti-starvation, in un
caso anche il vincolo sul campo visivo e il termine di ricompensa
dedicato all'attesa, degradano marcatamente su ogni metrica. Che la sola
assenza del termine di ricompensa non basti a spiegare questo degrado lo
conferma il confronto con la variante che rimuove solo quel termine
lasciando intatti stato e maschera: pur priva di qualunque penalità
sull'attesa nella ricompensa, resta nettamente sotto le due varianti
appena descritte sulle metriche di caso peggiore. Il comportamento
osservato dipende quindi più probabilmente dalla combinazione di più
meccanismi che da un singolo dettaglio implementativo, dell'algoritmo di
apprendimento o della sola ricompensa.

Il confronto tra meccanismi di comunicazione spaziale aggiunge una terza
componente. A parità di stato, ricompensa e maschera, il meccanismo che
assegna a ogni coppia di nodi vicini un peso di aggregazione fissato
dalla sola topologia, invece che appreso, resta sistematicamente indietro
su ogni metrica, in entrambi i regimi, arrivando a un caso peggiore
persino di Max-Pressure sul tempo di percorrenza massimo e sull'attesa
massima anche in test. Un secondo livello di
propagazione migliora questa variante su ogni metrica, pur restando
lontana dalle prestazioni delle altre due famiglie confrontate, mentre lo
stesso secondo livello peggiora il meccanismo di riferimento su ogni
metrica riportata, in entrambi i regimi. Tra le due profondità del
meccanismo di propagazione a onda, quella con più iterazioni risulta
leggermente migliore sul tempo di percorrenza massimo in entrambi i
regimi, coerente con il vantaggio teorico di un campo recettivo più
ampio; il modello di riferimento resta comunque quello con l'attesa
massima di test più bassa tra tutti i meccanismi confrontati. Equità ed
efficienza non sono quindi solo una funzione della ricompensa: richiedono
anche un meccanismo di comunicazione capace di instradare il segnale del
vicino in difficoltà fino a chi può agire su di esso.

Questi risultati confermano, entro i limiti discussi più sotto, i
contributi dichiarati nell'introduzione. L'ambiente di simulazione
costruito per questo lavoro ha permesso di condurre i tre studi appena
riportati sulle stesse configurazioni di training e su sette
configurazioni di generalizzazione; lo studio sul peso di equità ha
quantificato il compromesso con Max-Pressure senza ridurlo a un semplice
collasso verso il suo comportamento quando il peso viene azzerato; lo
studio di ablazione ha isolato singolarmente ogni modifica rispetto alla
procedura di riferimento, individuando nella combinazione di stato e
ricompensa, non in un singolo dettaglio dell'algoritmo di apprendimento,
la determinante principale del comportamento osservato; il confronto tra
meccanismi spaziali ha mostrato che il vantaggio osservato non è
specifico del meccanismo di attenzione appreso, ma neppure garantito da
una qualunque forma di condivisione spaziale dell'informazione; il
protocollo di valutazione a generalizzazione sistematica ha infine
permesso di verificare quali di questi risultati si mantengano, e in
quale misura, su condizioni mai incontrate in addestramento.

Nonostante i risultati ottenuti, alcuni limiti vanno considerati
nell'applicazione pratica di questo lavoro. Ogni modello, appreso o di
ablazione, deriva da un singolo run di addestramento
(§\ref{sec:protocollo}); la batteria a più seed introdotta nel protocollo
sperimentale ripete la sola valutazione, non l'addestramento, per cui
buona parte dei risultati più sottili discussi in questo capitolo resta
compatibile, oltre che con le spiegazioni proposte, con la varianza tra
run di addestramento indipendenti diversi dall'unico osservato. Il
confronto tra modelli in questo lavoro non include un test di
significatività statistica in senso classico: la consistenza osservata
su più configurazioni di test e più seed di valutazione rende i
risultati qualitativamente solidi, non un sostituto formale di un test
di ipotesi. La selezione del checkpoint migliore si basa sul solo tempo
di percorrenza medio (§\ref{sec:protocollo}), non sulle metriche di caso
peggiore che questo lavoro tiene in primo piano altrove, un limite
pratico del protocollo sperimentale più che una scelta di principio. La
valutazione di generalizzazione resta
confinata a reti sintetiche a griglia regolare, generate
parametricamente, mai a una rete stradale reale, un limite che questo
lavoro condivide con la maggior parte della letteratura citata
nell'introduzione, e che la sola estensione a griglie più grandi non
risolve. Allo stesso modo, \textsc{CityFlow} privilegia la velocità di
simulazione sul realismo microscopico del singolo veicolo
(§\ref{sec:background-tsc}), con alcune semplificazioni strutturali
dichiarate nella descrizione dell'ambiente
(§\ref{subsec:assunzioni-simulazione}): l'intervallo di giallo è
assimilato cinematicamente al verde, l'intervallo di tutto-rosso è
eseguito a ogni passo di decisione indipendentemente dal cambio di fase,
il limite di velocità è una costante uguale per ogni corsia e la
simulazione modella solo il traffico veicolare privato. Risultati
assoluti su un simulatore più dettagliato, o su dati reali, potrebbero
differire pur lasciando invariate le conclusioni relative tra i
controllori confrontati. Infine, tra i preset di ablazione quello che
rimuove il maggior numero di meccanismi modifica insieme più aspetti
dell'osservazione e della ricompensa, per cui i suoi risultati misurano
l'effetto congiunto di queste modifiche e non isolano un singolo
meccanismo (§\ref{sec:ablation-risultati}).

I risultati discussi finora aprono la strada a diverse direzioni di
ricerca. Una prima riguarda la gestione dell'interverde
(§\ref{subsec:interverde}): l'intervallo di tutto-rosso viene oggi
eseguito a ogni passo di decisione anche quando la fase successiva
coincide con quella corrente, per cui due dei quindici secondi di ogni
step sono tempo morto anche in assenza di un cambio di fase. Evitare il
tutto-rosso quando la stessa fase viene scelta in step consecutivi
manterrebbe il verde continuo in quel caso, restituendo tempo utile al
deflusso senza modificare l'intervallo di sicurezza nei casi in cui la
fase cambia davvero.

Una seconda direzione riguarda la maschera sulle azioni
(§\ref{sec:formalizzazione}). Il vincolo anti-starvation attuale vieta di
ripetere la stessa fase più di due volte consecutive, ma non limita per
quanto tempo complessivo una singola corsia possa restare rossa se le
fasi si alternano senza mai servirla. Una maschera per corsia, che ponga
un tetto al numero massimo di fasi consecutive attivate senza servire una
data corsia, porrebbe per costruzione un limite all'attesa massima,
agendo direttamente sulle metriche di caso peggiore che questo lavoro
tiene in primo piano; il costo in termini di efficienza di un vincolo di
questo tipo resta da studiare.

Le altre direzioni riguardano la generalizzazione a reti stradali
strutturalmente diverse da quelle qui testate. Le griglie di questo
lavoro hanno segmenti stradali di lunghezza uniforme, un vincolo che una
rete con strade di lunghezza eterogenea rimuoverebbe, mettendo alla prova
un modello il cui campo visivo e la cui nozione di distanza tra nodi sono
stati calibrati proprio su quella uniformità. Le intersezioni sono
inoltre tutte a quattro vie e ad angolo retto, mentre una griglia
irregolare, con angoli diversi da novanta gradi o incroci a tre o cinque
vie, metterebbe alla prova la struttura a fasi fisse qui adottata. Le vie
a senso unico cambierebbero a loro volta la struttura dei conflitti tra
movimenti alla base della stessa definizione di fase, un asse di
variazione distinto da quelli già esplorati. Più in generale, ciascuna
delle semplificazioni elencate tra le assunzioni del modello di
simulazione (§\ref{subsec:assunzioni-simulazione}) è un candidato
naturale per un futuro studio di generalizzazione, così come l'estensione
della valutazione a topologie stradali reali oltre alle griglie
sintetiche qui utilizzate.

Vale la pena di notare come molte di queste differenze strutturali, la
lunghezza di una strada, il suo limite di velocità, l'angolo di un
incrocio, potrebbero essere incorporate senza modificare l'architettura
di base, aggiungendole semplicemente alle meta-feature già usate dal
meta-learner di questo lavoro per condizionare i pesi di ciascun nodo
sulle sue caratteristiche locali. Un'ultima direzione, di natura diversa
dalle altre, è l'introduzione di mezzi pubblici e pedoni accanto al solo
traffico veicolare privato modellato in questo lavoro: un'estensione che
richiederebbe fasi e segnali dedicati, non la sola aggiunta di una nuova
meta-feature.

Alcune verifiche a basso costo permetterebbero infine di irrobustire i
risultati di questo lavoro senza modificarne l'impostazione. Ripetere
l'addestramento con più seed indipendenti per ciascun modello
confermerebbe se i risultati più sottili discussi in questo capitolo
riflettano davvero le modifiche studiate e non la sola variabilità tra
run di addestramento. Estendere inoltre la selezione del checkpoint
anche alle metriche di caso
peggiore, oggi basata sul solo tempo di percorrenza medio, verificherebbe
se questo cambio nella procedura di validazione porti a scegliere modelli
diversi da quelli selezionati in questo lavoro.
